\documentclass[11pt]{article}
\usepackage[margin=1in]{geometry}
\usepackage{booktabs}
\usepackage{longtable}
\usepackage{xcolor}
\usepackage{float}
\usepackage{hyperref}
\usepackage{enumitem}
\usepackage{fancyhdr}

\definecolor{fail}{RGB}{192,57,43}
\definecolor{flag}{RGB}{211,84,0}
\definecolor{pass}{RGB}{39,174,96}

\pagestyle{fancy}
\fancyhf{}
\rhead{\small Forensic Audit: Rajizadeh et al.\ (2017)}
\lhead{\small Paperscope}
\cfoot{\thepage}

\title{Forensic Statistics Audit\\[0.3em]
\large Rajizadeh et al., ``Effect of magnesium supplementation on depression status,''\\
\textit{Nutrition} 35 (2017) 56--60}
\author{Generated by Paperscope\\
\small \url{https://github.com/todd866/paperscope}}
\date{25 March 2026}

\begin{document}
\maketitle

\section*{What This Report Is}

This report checks whether the \textbf{numbers} in a published scientific paper are internally consistent. It does not evaluate whether the study was well-designed or whether magnesium actually helps depression. Instead, it asks a simpler question: \textit{could the reported statistics have come from real data?}

The answer, for this paper, is no --- not for all of them.

We used 19 automated checks based on techniques from the forensic metascience literature, primarily Heathers (2025).\footnote{Heathers, J. (2025). \textit{An Introduction to Forensic Metascience}. DOI: 10.5281/zenodo.14871843} These same techniques have been used to identify data problems in hundreds of published papers, leading to corrections and retractions.

\section*{What the Paper Claims}

Rajizadeh et al.\ conducted a small trial with 60 depressed patients who also had low magnesium. Half received magnesium supplements; half received a placebo. After 8 weeks, they measured depression scores and blood magnesium levels.

The paper concludes that magnesium supplementation improves both depression and magnesium levels. The trial was published in \textit{Nutrition}, a peer-reviewed Elsevier journal.

\section*{What We Found: The Summary}

We ran 50 individual checks across the paper's three data tables. The results:

\begin{center}
\begin{tabular}{lrl}
\toprule
\textbf{Category} & \textbf{Count} & \textbf{Meaning} \\
\midrule
\textcolor{fail}{\textbf{Failures}} & \textcolor{fail}{\textbf{14}} & Mathematically impossible \\
\textcolor{flag}{\textbf{Flags}} & \textcolor{flag}{\textbf{7}} & Suspicious but not proven impossible \\
\textcolor{pass}{Passes} & \textcolor{pass}{29} & Checks out fine \\
\midrule
\textbf{Total} & \textbf{50} & \\
\bottomrule
\end{tabular}
\end{center}

\noindent Verdict: \textbf{severe data integrity concerns}. Multiple reported statistics cannot have come from the data described.

\section{The Checks, Explained}

\subsection{GRIM: Can This Average Be Real?}

\textbf{The idea:} The Beck Depression Inventory (BDI) is a questionnaire where patients score each of 21 items as 0, 1, 2, or 3. The total is always a whole number between 0 and 63. When you average whole numbers from a known number of people, only certain averages are mathematically possible.

\textbf{Example:} If 26 people each score a whole number, the group average must be some whole number divided by 26. The possible averages near 11 are $284/26 = 10.923$, $285/26 = 10.961$, $286/26 = 11.000$, $287/26 = 11.038$, and so on. The reported average of 11.26 would require $26 \times 11.26 = 292.76$ --- but you can't have 0.76 of a point. No combination of 26 whole numbers produces an average of exactly 11.26.

\textbf{Result:}

\begin{center}
\begin{tabular}{llrl}
\toprule
\textbf{Group} & \textbf{Timepoint} & \textbf{Mean} & \textbf{Verdict} \\
\midrule
Magnesium & Baseline & 26.9 & \textcolor{pass}{Pass} \\
Magnesium & End & 11.26 & \textcolor{fail}{\textbf{Impossible}} \\
Placebo & Baseline & 25.6 & \textcolor{pass}{Pass} \\
Placebo & End & 15.2 & \textcolor{pass}{Pass} \\
\bottomrule
\end{tabular}
\end{center}

\subsection{SPRITE: Can Any Dataset Produce These Statistics?}

\textbf{The idea:} GRIM checks the average alone. SPRITE goes further: given the reported average \textit{and} the spread (standard deviation), is there \textit{any} possible set of whole-number BDI scores that produces both? The algorithm randomly shuffles candidate datasets and tries to match both statistics simultaneously.

\textbf{Result:} Using 50,000 iterations per seed with 3 independent seeds per target, three of the four BDI mean/SD pairs could not be reconstructed. One (placebo end) was found in 515 attempts, confirming it is achievable. A separate deep search using 2 million iterations per seed confirmed the same pattern, with the closest near-misses still falling short.

\subsection{Correlation Bound: Are the Spreads Mutually Consistent?}

\textbf{The idea:} When you measure something before and after treatment, you get three spreads: the spread of ``before'' scores, the spread of ``after'' scores, and the spread of the \textit{changes}. These three numbers are mathematically linked by a formula involving the correlation between before and after. If you plug in the reported numbers and the formula says the correlation must be greater than 1 (or less than $-1$), the numbers are impossible --- correlations are bounded between $-1$ and $+1$.

\textbf{Result:} For the placebo group's serum magnesium measurements:
\begin{itemize}[nosep]
\item Before SD = 0.13, After SD = 0.27, Change SD = 0.03
\item Implied correlation: $r = 1.27$ --- \textcolor{fail}{\textbf{impossible}}.
\item The smallest possible change SD for these before/after SDs is 0.14, not 0.03.
\end{itemize}

This is one of the strongest findings. The reported change SD of 0.03 is off by nearly a factor of 5 from the minimum that mathematics allows.

\subsection{P-Value Recalculation}

\textbf{The idea:} A p-value is calculated from the average, the spread, and the sample size using a specific formula. We can redo the calculation and check whether the reported p-value matches.

\textbf{Result:} The \textit{between-group} comparisons (the main results) check out --- the reported p-values match our calculations. But the \textit{within-group} paired tests are off by orders of magnitude:

\begin{center}
\begin{tabular}{lrrl}
\toprule
\textbf{Test} & \textbf{Reported p} & \textbf{Calculated p} & \textbf{Verdict} \\
\midrule
Placebo serum Mg change & 0.110 & $1.0 \times 10^{-14}$ & \textcolor{flag}{\textbf{Flag}} \\
Mg group serum Mg change & 0.001 & $1.2 \times 10^{-5}$ & \textcolor{flag}{\textbf{Flag}} \\
Mg group BDI change & 0.001 & $2.9 \times 10^{-9}$ & \textcolor{flag}{\textbf{Flag}} \\
Placebo BDI change & 0.001 & $2.9 \times 10^{-7}$ & \textcolor{flag}{\textbf{Flag}} \\
\bottomrule
\end{tabular}
\end{center}

\noindent The placebo serum Mg result is the most striking: the reported p = 0.110 (not significant) should be $p \approx 10^{-14}$ (astronomically significant) given the reported numbers. This is likely related to the impossible change SD of 0.03 discussed above.

\subsection{Arithmetic Consistency: Do the Tables Add Up?}

\textbf{The idea:} If baseline is 63.33 and endpoint is 71.51, then the change should be $71.51 - 63.33 = 8.18$. We check every row.

\textbf{Result:} Six of eight dietary intake rows in Table 2 fail this check, with discrepancies ranging from 28 to 169 units.

One hypothesis (proposed by Meyerowitz-Katz on PubPeer) is that the mean and standard deviation columns were accidentally swapped. This is plausible because, for example, in the placebo protein row the actual change is $69.67 - 75.08 = -5.41$, and the number reported in the ``SD'' position is $-5.40$ --- matching the true change, not a spread measure. The number in the ``mean change'' position (23.45) would then be the actual SD.

Additionally, three reported standard deviations in Table 2 are \textit{negative} ($-5.40$, $-55.89$, $-5.56$). Standard deviations measure spread and cannot be negative. If these are actually mean changes that ended up in the wrong column, the sign makes sense. However, we note this is an interpretation of the pattern, not a proven fact --- other explanations (e.g., data entry errors from a different source) are possible.

\subsection{Carlisle Test: Was the Randomization Real?}

\textbf{The idea:} In a properly randomized trial, the baseline characteristics of the two groups should differ by chance. The p-values comparing groups at baseline should be uniformly distributed. If they are \textit{too} perfectly balanced, the randomization may have been manipulated.

\textbf{Result:} The Table 1 baseline p-values (0.93, 0.28, 0.80, 0.67, 0.89) pass the Carlisle test --- they are consistent with real randomization. This is actually reassuring: the data integrity problems appear to be in the \textit{outcome} reporting, not the trial design.

\subsection{Other Checks}

\begin{itemize}[nosep]
\item \textbf{GRIMMER} (SD consistency for integer data): the Mg group endpoint fails because the mean already fails GRIM.
\item \textbf{Chi-squared recalculation} on Table 1: sex and marital status p-values match.
\item \textbf{SD/SE confusion}: the placebo serum Mg change SD of 0.03 is flagged as possibly being a standard \textit{error} rather than a standard deviation --- if it were the SE, the true SD would be $\sim$0.16, which is within plausible range.
\item \textbf{Effect size consistency}: the between-group BDI effect (Cohen's $d = -0.62$, $p = 0.028$) is internally consistent.
\item \textbf{Variance ratios}: group variances are not suspiciously similar or different.
\item \textbf{Quick SD check}: all BDI standard deviations are plausible for the 0--63 range.
\end{itemize}

\section{What Does This Mean?}

The pattern is clear:

\begin{enumerate}[nosep]
\item \textbf{Table 1} (baseline characteristics) and \textbf{between-group comparisons} are internally consistent. The randomization looks real, and the main treatment effect statistics check out.
\item \textbf{Table 2} (dietary intake changes) has systematic errors consistent with a mean/SD column swap during manuscript preparation.
\item \textbf{Table 3} (outcomes) contains at least one impossible mean, an impossible correlation bound, and p-values that don't match the reported statistics.
\end{enumerate}

These findings do \textbf{not} prove fraud. Possible explanations include transcription errors, software misuse, rounding cascades, or subset confusion (e.g., running an analysis on 25 people but reporting $n = 26$). However, they make it impossible to evaluate the paper's claims from the published numbers alone.

The appropriate response, as recommended by forensic metascience practice, is to request the raw dataset from the authors so these discrepancies can be resolved.

\section{How This Analysis Was Produced}

This report was generated using \textbf{Paperscope}, an open-source toolkit for academic paper analysis. The forensic statistics module implements 19 checks based on techniques described in Heathers (2025), Brown \& Heathers (2017), and related work.

\textbf{Important methodological note:} The values tested in this audit were \textit{manually transcribed} from the published tables in Rajizadeh et al.\ (2017), not automatically extracted from the PDF. This paper serves as the built-in demonstration case for the forensic module. The checks themselves are automated --- once the numbers are entered, no subjective judgment affects the results --- but the initial data entry step is manual and therefore subject to transcription error. All values can be verified against the original paper (Tables 1--3, pages 58--59).

In a typical review workflow, the reviewer reads the paper, transcribes the relevant summary statistics, and feeds them to the forensic functions. Automated PDF table extraction is a separate capability not used here.

\medskip
\noindent\textbf{Software:} \url{https://github.com/todd866/paperscope}\\
\textbf{Module:} \texttt{paperscope.analysis.forensic\_stats}\\
\textbf{Reference:} Heathers, J. (2025). \textit{An Introduction to Forensic Metascience}. DOI: 10.5281/zenodo.14871843

\newpage
\appendix
\section{Provenance Table}

Every tested value is listed below with its source location in the original paper, the check applied, and the result.

\small
\begin{longtable}{p{3.5cm}p{2.2cm}p{1.8cm}p{2.5cm}p{2.5cm}}
\toprule
\textbf{Source} & \textbf{Value} & \textbf{Check} & \textbf{Expected} & \textbf{Result} \\
\midrule
\endfirsthead
\toprule
\textbf{Source} & \textbf{Value} & \textbf{Check} & \textbf{Expected} & \textbf{Result} \\
\midrule
\endhead

\multicolumn{5}{l}{\textit{Table 3, Beck score means (integer instrument, BDI range 0--63)}} \\
\midrule
T3, Mg baseline & mean=26.9, n=26 & GRIM & n$\times$mean $\in \mathbb{Z}$ & \textcolor{pass}{Pass} \\
T3, Mg end & mean=11.26, n=26 & GRIM & n$\times$mean $\in \mathbb{Z}$ & \textcolor{fail}{Fail: 292.76} \\
T3, Placebo baseline & mean=25.6, n=27 & GRIM & n$\times$mean $\in \mathbb{Z}$ & \textcolor{pass}{Pass} \\
T3, Placebo end & mean=15.2, n=27 & GRIM & n$\times$mean $\in \mathbb{Z}$ & \textcolor{pass}{Pass} \\
\midrule

\multicolumn{5}{l}{\textit{Table 3, Beck score mean+SD pairs (SPRITE, 50k iter $\times$ 3 seeds)}} \\
\midrule
T3, Mg baseline & 26.9 $\pm$ 7.1, n=26 & SPRITE & any valid dataset & \textcolor{flag}{Not found} \\
T3, Mg end & 11.26 $\pm$ 6.9, n=26 & SPRITE & any valid dataset & \textcolor{fail}{Fail (GRIM)} \\
T3, Placebo baseline & 25.6 $\pm$ 6.1, n=27 & SPRITE & any valid dataset & \textcolor{flag}{Not found} \\
T3, Placebo end & 15.2 $\pm$ 9.3, n=27 & SPRITE & any valid dataset & \textcolor{pass}{Found (iter 515)} \\
\midrule

\multicolumn{5}{l}{\textit{Table 3, Serum Mg pre/post/change SDs (correlation bound)}} \\
\midrule
T3, Placebo Mg & SD: 0.13, 0.27, 0.03 & Corr.\ bound & $|r| \leq 1$ & \textcolor{fail}{Fail: $r$=1.27} \\
T3, Mg group & SD: 0.19, 0.19, 0.29 & Corr.\ bound & $|r| \leq 1$ & \textcolor{pass}{Pass: $r$=$-$0.16} \\
\midrule

\multicolumn{5}{l}{\textit{Table 3, Paired t-test p-values}} \\
\midrule
T3, Placebo Mg change & p=0.110 & Recalc & match & \textcolor{flag}{$p_\text{calc}$=$10^{-14}$} \\
T3, Mg group Mg change & p=0.001 & Recalc & match & \textcolor{flag}{$p_\text{calc}$=$10^{-5}$} \\
T3, Mg BDI change & p=0.001 & Recalc & match & \textcolor{flag}{$p_\text{calc}$=$10^{-9}$} \\
T3, Placebo BDI change & p=0.001 & Recalc & match & \textcolor{flag}{$p_\text{calc}$=$10^{-7}$} \\
\midrule

\multicolumn{5}{l}{\textit{Table 3, Between-group t-tests}} \\
\midrule
T3, BDI baseline & p=0.49 & Recalc & match & \textcolor{pass}{Pass: 0.48} \\
T3, BDI end & p=0.08 & Recalc & match & \textcolor{pass}{Pass: 0.09} \\
T3, BDI change & p=0.02 & Recalc & match & \textcolor{pass}{Pass: 0.03} \\
T3, Mg baseline & p=0.27 & Recalc & match & \textcolor{pass}{Pass: 0.27} \\
\midrule

\multicolumn{5}{l}{\textit{Table 2, Change scores (End $-$ Baseline = Change?)}} \\
\midrule
T2, Protein, Mg & change=44.31 & Arithmetic & 8.18 & \textcolor{fail}{Fail ($\Delta$36)} \\
T2, Protein, Placebo & change=23.45 & Arithmetic & $-$5.41 & \textcolor{fail}{Fail ($\Delta$29)} \\
T2, Carbs, Mg & change=87.21 & Arithmetic & 2.39 & \textcolor{fail}{Fail ($\Delta$85)} \\
T2, Carbs, Placebo & change=113.34 & Arithmetic & $-$55.89 & \textcolor{fail}{Fail ($\Delta$169)} \\
T2, Fat, Mg & change=$-$6.35 & Arithmetic & $-$6.35 & \textcolor{pass}{Pass} \\
T2, Fat, Placebo & change=32.15 & Arithmetic & $-$5.57 & \textcolor{fail}{Fail ($\Delta$38)} \\
T2, Diet.\ Mg, Mg & change=81.50 & Arithmetic & 17.36 & \textcolor{fail}{Fail ($\Delta$64)} \\
T2, Diet.\ Mg, Placebo & change=90.29 & Arithmetic & 30.14 & \textcolor{fail}{Fail ($\Delta$60)} \\
\midrule

\multicolumn{5}{l}{\textit{Table 2, Negative SDs}} \\
\midrule
T2, Protein $\Delta$, Placebo & SD=$-$5.40 & SD $\geq$ 0 & non-negative & \textcolor{fail}{Fail} \\
T2, Carbs $\Delta$, Placebo & SD=$-$55.89 & SD $\geq$ 0 & non-negative & \textcolor{fail}{Fail} \\
T2, Fat $\Delta$, Placebo & SD=$-$5.56 & SD $\geq$ 0 & non-negative & \textcolor{fail}{Fail} \\
\midrule

\multicolumn{5}{l}{\textit{Table 1, Baseline p-values (Carlisle test + chi-squared recalc)}} \\
\midrule
T1, 5 baseline p-values & 0.93--0.89 & Carlisle & uniform on [0,1] & \textcolor{pass}{Pass: Z=$-$1.52} \\
T1, Sex & $\chi^2$, p=0.93 & Chi-sq recalc & match & \textcolor{pass}{Pass: p=0.93} \\
T1, Marital status & $\chi^2$, p=0.28 & Chi-sq recalc & match & \textcolor{pass}{Pass: p=0.29} \\

\bottomrule
\end{longtable}
\normalsize

\end{document}
